%\section{Introduction}

The recent discovery that the neutrino mixing angle $\theta_{13}\approx 9^\circ$~\cite{Abe:2011sj,Adamson:2013ue,Abe:2011fz,An:2012eh,Ahn:2012nd} makes measuring the hierarchy of neutrino masses and CP violation possible in precision neutrino oscillation experiments.  
Quasi-elastic interactions, $\overline{\nu} p \rightarrow \ell^+  n$ and  $\nu n \rightarrow \ell^- p$, have simple kinematics and serve as reference processes in those experiments ~\cite{AguilarArevalo:2008rc,Abe:2011sj,Ayres:2004js} at GeV energies. These processes are typically modeled as scattering on free nucleons in a relativistic Fermi gas (RFG), with a nucleon axial form factor measured in neutrino-deuterium quasi-elastic scattering~\cite{Bodek:2007vi,Kuzmin:2007kr}. In the RFG model~\cite{Smith:1972xh} the initial state nucleons are independent in the mean field of the nucleus, and therefore the neutrino energy and momentum transfer $Q^2$ can be estimated from the polar angle $\theta_\ell$ and momentum $p_\ell$ of the final state lepton. 
However, correlations and motion of the initial state nucleons, as well as interactions of the final state particles within the nucleus, significantly modify the Fermi gas picture and affect the neutrino energy reconstruction in oscillation experiments~\cite{Martini:2012uc,Lalakulich:2012hs,Nieves:2012yz}. 

Few measurements of antineutrino quasi-elastic scattering exist~\cite{Bonetti:1977cs,Ahrens:1988rr,AguilarArevalo:2013hm}.  
The most recent, from the MiniBooNE experiment on a hydrocarbon target at energies near $1$~GeV~\cite{AguilarArevalo:2013hm}, does not agree with expectations based on the RFG model described above.  A MiniBooNE analysis of $\nu_\mu$ quasi-elastic scattering suggests an increased axial form factor at high $Q^2$~\cite{AguilarArevalo:2007ab}. 
However, results at higher energy from the NOMAD experiment~\cite{Lyubushkin:2008pe} are consistent with the Fermi gas model and the form factor from deuterium.

In this Letter we report the first study of antineutrino quasi-elastic interactions from the \minerva experiment, which uses a finely segmented scintillator detector at Fermilab to measure muon antineutrino and neutrino charged current interactions at energies between $1.5$ and $10$~GeV on nuclear targets.  
The signal reaction has a $\mu^+$ in the final state along with one or more nucleons (typically with a leading neutron), and no mesons\footnote{{I}n this analysis quasi-elastic scattering occurs on both free protons and inside carbon nuclei.}. 
The $\mu^+$ is identified by a minimum ionizing track that traverses  \minerva~\cite{minerva_nim} and travels downstream to the MINOS magnetized spectrometer~\cite{Michael:2008bc} where its momentum and charge are measured. 
The leading neutron, if it interacts, leaves only a fraction of its energy in the detector in the form of scattered low energy protons.  
To isolate quasi-elastic events from those where mesons are produced, we require the hadronic system recoiling against the muon to have a low energy. That energy is measured in two spatial regions. The {\it vertex energy} region corresponds to a sphere around the vertex with a radius sufficient to contain a proton (pion) with 120 (65)~MeV kinetic energy. This region is sensitive to low energy protons which could arise from correlations among nucleons in the initial state or interactions of the outgoing hadrons inside the target nucleus. We do not use the vertex energy in the event selection.  The {\it recoil energy} region includes energy depositions outside of the vertex region and is sensitive to pions and higher energy nucleons. We use the recoil energy to estimate and remove inelastic backgrounds.




